\chapter{Testing and validation}
\label{ch:testing}

Yap has no client test suite. Correctness is established by static checks and continuous integration on
every change, a small set of SQL tests for the server rules, behaviour proofs run against a disposable
backend for changes that alter behaviour, and manual testing on phones. This chapter describes
each and states plainly what is not covered. \Cref{fig:validation} shows the path a change takes.

\begin{figure}[ht]
  \centering
  % How a change is validated before it reaches main (Figure: validation).
\begin{tikzpicture}[
    stage/.style={box, font=\sffamily\scriptsize, text width=28mm, minimum height=24mm, align=flush left, anchor=north},
    lab/.style={font=\sffamily\scriptsize, text=muted, fill=white, inner sep=1.2pt},
  ]
  \node[stage, clientbox, text width=20mm, align=center] (change) at (1.2,0) {\textbf{Change}\\[3pt]on a branch, with its intent};
  \node[stage] (static) at (4.35,0) {\textbf{Static checks}\\[3pt]\texttt{npm run typecheck}\\\texttt{expo lint}\\\texttt{expo-doctor}\\[2pt]{\color{muted}larger changes add \texttt{expo export}, iOS and Android}};
  \node[stage, serverbox] (proofs) at (7.95,0) {\textbf{Behaviour proofs}\\[3pt]local Supabase stack\\before the change: fails\\after the change: passes};
  \node[stage, accentbox] (pipeline) at (11.55,0) {\textbf{Review and CI}\\[3pt]review against the intent; pull request; GitHub Actions: app checks, migration files, database rebuild and tests};
  \node[stage, clientbox, text width=14mm, align=center] (main) at (14.95,0) {\textbf{main}};

  \draw[flow] (change.east |- 0,-1.2) -- (static.west |- 0,-1.2);
  \draw[flow] (static.east |- 0,-1.2) -- (proofs.west |- 0,-1.2);
  \draw[flow] (proofs.east |- 0,-1.2) -- (pipeline.west |- 0,-1.2);
  \draw[flow] (pipeline.east |- 0,-1.2) -- (main.west |- 0,-1.2);
  \node[lab, anchor=north, text width=32mm, align=center] at ($(proofs.south)+(0,-0.1)$) {for changes to behaviour or to the database};

  \node[box, font=\sffamily\scriptsize, text width=58mm, align=center, anchor=north] (manual) at (4.35,-3.3)
    {\textbf{Manual testing}: two phones in Expo Go, two accounts, close together, both live};
  \draw[softflow] (change.south) |- (manual.west);
\end{tikzpicture}

  \caption{Validation of a change before it reaches \code{main}.}
  \label{fig:validation}
\end{figure}

\section{Static checks}

The README lists three checks, and continuous integration adds a fourth:

\begin{description}
  \item[Type check.] \code{npm run typecheck} type-checks the whole client in TypeScript's strict mode,
    including the typed routes and the database types generated from the schema, so a call with the wrong
    function arguments or a link to a missing route fails to compile. The route types are generated, not
    committed, so the script first writes them to \code{.expo/types} with the command
    \code{expo customize tsconfig.json} and only then runs \code{tsc --noEmit}. Without that step a fresh
    checkout, such as the one in CI, would accept any route string.
  \item[Lint.] \code{npm run lint} runs ESLint with Expo's configuration and the rules of the ESLint
    security plugin, which flag risky patterns such as \code{eval}. A web rule about apostrophes in markup is turned off because it
    does not apply to React Native text, and the security rule that flags every \code{object[key]} is
    turned off because TypeScript already checks those.
  \item[Project health.] \code{npx expo-doctor} checks that the installed packages match the Expo SDK.
  \item[Release bundles.] Exporting with \code{npx expo export}, once with \code{--platform ios} and once
    with \code{--platform android}, builds the production JavaScript bundles and so proves that the app
    bundles for both platforms.
\end{description}

\section{Continuous integration}
\label{sec:ci}

GitHub Actions (\code{.github/workflows/ci.yml}) runs on every pull request and every push to
\code{main}, and a newer push to the same branch cancels the older run. Three jobs report separately:

\begin{description}
  \item[App.] \code{npm ci}, the type check, the lint with warnings treated as failures, \code{expo-doctor}
    and both release bundles, plus an \code{npm audit} report of the production dependencies that never
    fails the run.
  \item[Migration files.] Every migration is named \code{<14-digit version>\_<name>.sql}; a new one must be
    later than every migration on the target branch; and an existing migration that is edited or deleted
    fails the check (\code{scripts/check-migrations.sh}). A rename that leaves the file byte for byte the
    same passes, so a file can take the version the hosted project recorded (\cref{app:migrations}).
  \item[Database (reset, lint, tests, types).] Starts Supabase in Docker with a pinned CLI version, applies
    every migration to an empty database with \code{supabase db reset}, lints the functions in
    \code{public} and \code{private} with \code{supabase db lint}, runs the SQL tests in
    \code{supabase/tests} with \code{supabase test db}, and fails when \code{src/data/database.types.ts}
    differs from the types generated from the migrations (\code{npm run db:types:check}).
\end{description}

Dependabot opens weekly update pull requests for npm packages and GitHub Actions; packages whose version
the Expo SDK decides only move in an SDK upgrade. An optional AI review (PR-Agent, with a Gemini API key)
stays off until the owner turns it on. \code{docs/RELEASING.md} describes the whole flow around these
checks: branches and squash-merged pull requests, the branch protection that should require the three
jobs (not yet switched on), versioning, EAS builds and store submission, applying database changes,
rolling back, dependency updates and the optional review (\cref{sec:release}).

\section{SQL tests}

The pgTAP file \code{supabase/tests/hardening_test.sql} runs against the local stack with
\code{supabase test db}, inside a transaction it rolls back. Its 19 checks cover what a modified app could do
before the October 2026 hardening migration: names, intros, interests and messages that show nothing
or carry direction overrides or stacked accents, a second say-hi creating a second request, repeated
reports, deleting events to get around the daily limit, uploads into a sub-folder or past 16 files, and
\code{my_status} without the server clock. One more checks that the role the dashboard uses,
\code{service_role}, can still edit a profile now that its check constraints call private helpers.

The second file, \code{supabase/tests/parallel_limits_test.sql}, sends calls at the same moment from separate connections (through \code{dblink}): ten events, twelve people
reported and five reports of one person. Its three checks pass only if the daily limits hold, five
events, ten reports and one report. Before \code{create_event} and \code{report_user} took a per-person
lock, all three failed (ten events, twelve reports, five reports). The people it signs up through those
connections are saved for real, so it deletes them at its start and end. Continuous integration runs
every file on every pull request, in the Database job. The photo limits of 16 files and 40 uploads a day
have no such check: they are best-effort against uploads sent at once.

The third file, \code{supabase/tests/safety_test.sql}, has 51 checks for the October 2026 safety
migration, each something the QA report did before it: a live position is stored as a cell centre and the
accuracy is clamped; changing only the accuracy never changes who is nearby; \code{nearby} and
\code{events_near} refuse the thirteenth and twenty-first call in a minute; three brand-new accounts that
only said hi, shared an event or never met someone cannot hide them, while a say-hi from the person, a
chat or meeting with them, a chat they accepted and then ended, or their message in a shared group chat
makes a report count (a say-hi taken back before an answer does not), three such reports hide them, and
the hide ends the person's live session, waiting request and event membership;
a person restored after review is a member of the event they host again and is not hidden again by one
more report; a hidden account is refused with \code{account restricted}
when it says hi, sends either kind of message, changes its note or goes live; the sign-up hook refuses a
suspended person's address with a \code{+tag}, a throwaway subdomain and a twenty-first sign-up from one
IP address, and \code{restore_user} lifts the ban; deleting a hidden account bans its email and keeps the
seven reports about it; the map lists only events around the position the server holds, hides the host's
photo until the user joins, refuses a jump to another city (also by going offline and live again, and
after the map position is older than two hours) and a fourth location change in ten minutes, accepts a
commute, lets going live replace an older map position, and \code{join_event} measures from the held
position; the bucket limit is 1.5\,MiB and a forty-first
upload in a day is refused; and an event category outside the app's twelve is refused.

The fourth file, \code{supabase/tests/restricted_visibility_test.sql}, has 34 checks for the
\code{restricted_visibility} migration. For a person hidden after reports and for one suspended by a
moderator, each of whom met one person who said hi to them and one they said hi to, it checks that both
partners lose the card from \code{request_person} and can read none of the profile picture, its small copy
or the gallery photo; that the restricted person still reads their own photos but gets no card or photos
of the people they met; and that \code{restore_user} makes the card and photos visible to both partners again.

\section{Pre-merge review pipeline}

Recent changes reach \code{main} through a pre-push pipeline, \code{no-mistakes}, which reviews the change
against its stated intent, runs the checks above, checks that documentation still matches the code, and
opens the pull request. Each finding is either fixed in a separate commit or recorded in the pull request,
so the history of a change shows what was found and how it was resolved.

\section{Behaviour proofs against a disposable backend}

Changes that alter behaviour were proved against a throwaway local Supabase stack running in Docker, built
from the repository's migrations, with real signed-up users. Each proof ran the scenario twice: on the code
before the change, where it was expected to fail, and on the change, where it had to pass. The hosted
project and \code{src/config/supabase.ts} were never touched; the client code under test was pointed at
the local stack through a module alias. Two pull requests show the method.

\paragraph{Client reliability (pull request 61)} The app's real store and API layer were bundled for
Node, with small stand-ins for React Native, the network module and AsyncStorage, and driven as one phone
while a plain Supabase client acted as the other. Twelve scenarios covered missed chat endings, requests
lost in the background, two chats at once, sign-out order, a failed profile load, repeated photo
downloads, raw database errors, a failed go-offline and the history limit. All twelve reproduced the bug
before the change and passed after it. One item, leaving a group chat before its history loads, had no
recorded proof.

\paragraph{Server rules (pull request 63)} A suite of SQL-level scenarios made the same function, table
and storage calls as the app and covered every new rule: cancelling, one chat at a time, blocks, photo
paths and photo visibility, relocation limits, stale chats, the event radius and join distance, the new
return values and the function grants. On the previous migrations it recorded 51 passes and 56 failures;
with the new migration, 118 passes and no failures. A second suite ran the real store and API layer as
four simulated phones over Realtime and passed all 21 of its steps, from the server-time countdown to
blocking during a chat; on the previous code it failed at the second and sixth steps and stopped at the seventh. The
generated database types were regenerated from the local stack as part of the same change.

\paragraph{Trying to break the app (October 2026)} Four users signed up on the local stack and made 83
attempts directly over HTTP, as a modified app would, against the tables, functions and storage paths the
app uses: extreme text in every field, double taps on every action, other users' IDs, out-of-range values
and storage abuse. On the previous migrations 28 abusive attempts got through (26 accepted outright, plus
twelve copies of one report and a report pointing at a stranger's request), and a second say-hi answered
with an error. With the hardening migration all 28 are refused or made harmless, a second say-hi returns
the first request, and the 29 attempts the earlier rules refused are still refused. The SQL tests above
came out of this pass: on the previous migrations the first file fails its first five checks and stops
at the eighth, and the second fails all three; with the migration all 22 pass. The generated database types came out unchanged, because no public
function changed its arguments or return type.

\section{Manual testing}

The core loop is tested by hand with two phones running Expo Go, signed in with different email
addresses, close to each other (the same room always counts) and both live. The README describes this setup and the
dashboard settings it needs.

\section{What is not covered}

\begin{itemize}
  \item There are no unit tests, including for the pure helpers in \code{src/lib}, which would be the
    easiest to test.
  \item There are no automated user-interface tests on a device or simulator.
  \item The behaviour-proof scripts were written for their pull requests and are not part of the
    repository, so they cannot be rerun as a regression suite; only \code{supabase/tests} is, and
    continuous integration runs it on every change.
  \item The word filter has no SQL test in \code{supabase/tests}, and \code{private.suspend_user} is tested
    only for the ban it leaves.
\end{itemize}
